\documentclass[11pt,a4paper]{article}

\usepackage[T1]{fontenc}
\usepackage[utf8]{inputenc}
\usepackage{lmodern}
\usepackage[margin=2.5cm]{geometry}
\usepackage{microtype}
\usepackage{enumitem}
\usepackage{xcolor}
\usepackage{hyperref}

\setlength{\parindent}{0pt}
\setlength{\parskip}{0.6em}
\setlist[itemize]{topsep=0.2em,itemsep=0.25em}

\definecolor{linkblue}{HTML}{1F4E79}
\hypersetup{
    colorlinks=true,
    linkcolor=linkblue,
    urlcolor=linkblue,
    pdftitle={VIP Week 9 Report},
    pdfauthor={Yorgo Azzi}
}

\begin{document}

\begin{center}
    {\Large\bfseries VIP Week 9 Report}\\[0.5em]
    {\large Yorgo Azzi}\\[0.25em]
    November 11, 2026
\end{center}

\section{Overview}

This week, I developed a new Conductor mode for the movement-controlled music
system. The objective was to give the performer direct control over when the
music advances, instead of allowing the model to continue according to a fixed
clock. MediaPipe hand tracking remains the input method, while the two hands
now have separate and clearly defined musical roles.

The right hand controls the beat and chord selection. The left hand performs
the soprano melody through vertical movement. The Markov model continues to
generate the three accompaniment voices and maintain the musical structure.

\section{Work Completed}

\subsection{Event-Driven Conductor Mode}

A third interface mode, called \textbf{Conductor}, was added alongside Full
SATB and Soprano Performance. In this mode, the accompaniment does not advance
automatically. It waits for a conducting gesture from the performer.

MediaPipe provides the position of the hand landmarks in each camera frame. A
stable palm center is calculated from the wrist and knuckle landmarks. The
system then estimates palm velocity and acceleration over time. An acceleration
peak is treated as a beat command, and the time between accepted beats provides
an estimate of the performance tempo.

The detector was refined so that only a \textbf{downward right-hand stroke}
can trigger a beat. Upward recovery, sideways movement, and all left-hand
acceleration are ignored. I increased the acceleration threshold to 10 in the
current setup so that small tracking changes and ordinary hand motion do not
advance the music. A short refractory period also prevents one stroke from
being detected more than once.

When a beat is accepted, the realtime engine advances exactly one beat and
forces an audible onset. If no new downstroke is detected, the current notes
remain sustained. This makes the generation follow the performer's conducting
gestures rather than the internal clock.

\subsection{Right-Hand Chord Control}

The right hand also selects the chord degree that will be applied at the next
conducted beat. The finger combinations are:

\begin{itemize}
    \item index: first degree;
    \item index and middle: second degree;
    \item index, middle, and ring: third degree;
    \item index, middle, ring, and pinky: fourth degree;
    \item all five fingers: fifth degree;
    \item thumb: sixth degree;
    \item thumb and index: seventh degree.
\end{itemize}

The selected degree is interpreted within the chosen key and mode. The chord
gesture is sampled when the right-hand beat occurs, so the performer prepares
a finger pattern and then uses the downstroke to apply it. If the current
finger pattern is not mapped, the previous chord degree is retained.

An earlier version divided chord control between both hands, with degrees six
and seven assigned to the left hand. This was removed so that the right hand
has one complete role: controlling the timing and harmony.

\subsection{Left-Hand Melody Control}

The left hand now controls only the soprano melody. At every conducted beat,
its current height becomes the reference position. Until the following beat,
moving the hand upward produces successive ascending scale notes, while moving
it downward produces successive descending scale notes. A stationary hand
holds the current soprano note.

Only one scale step is emitted per camera frame. This allows a larger hand
movement to be heard as a short ascending or descending scale instead of
replacing several notes at the same instant. Right-hand height has no effect
on the melody, and left-hand movement cannot trigger the accompaniment.

\subsection{Realtime Integration}

Conductor mode uses an event queue between the camera interface and the audio
thread. Beat events advance the Markov accompaniment, while melody events
replace only the soprano note and sustain the other three voices. This keeps
melodic movement responsive without waiting for another accompaniment beat.

The mode was added to the existing interface with its own button and keyboard
shortcut. Its soprano output appears in the integrated piano roll, while Full
SATB and Soprano Performance remain available without changing their existing
controls.

\section{Results}

The main results were:

\begin{itemize}
    \item Music generation can now wait for an intentional conducting gesture
    instead of progressing continuously.
    \item Downward-only detection gives the right-hand stroke a clearer musical
    meaning and prevents the upward recovery from creating an extra beat.
    \item Raising the acceleration threshold reduces accidental beat commands
    caused by small movements and landmark jitter.
    \item The right hand provides a complete beat-and-chord control role,
    including all seven chord degrees.
    \item The left hand independently performs successive soprano scale notes,
    while the Markov model supplies the three accompaniment voices.
    \item The separation between the two hands makes the interaction easier to
    understand: right for timing and harmony, left for melody.
    \item The new mode works inside the existing camera, MIDI, piano-roll, and
    mode-selection interface.
\end{itemize}

\section{Challenges}

Acceleration estimates depend on camera frame timing and MediaPipe landmark
stability. A low threshold can create false beats, while a very high threshold
can require an exaggerated gesture or miss softer conducting strokes. The
value of 10 works as a stricter starting point for the current setup, but it
still requires physical testing under different camera positions and lighting
conditions.

The local acceleration maximum is confirmed one camera frame after it occurs.
This improves resistance to landmark noise but introduces a small detection
delay. The complete perceived response also includes camera capture, hand
tracking, event processing, MIDI output, and synthesizer latency.

Chord accuracy still depends on MediaPipe correctly recognizing the finger
pattern at the exact time of the beat. Thumb-based gestures for the sixth and
seventh degrees may require additional tuning because thumb extension is more
difficult to classify consistently than the other fingers.

Finally, the left-hand melody is limited to notes from the selected scale. This
keeps the result coherent with the accompaniment, but the performer currently
has no direct control over larger melodic intervals or notes outside the scale.

\section{Next Steps}

\begin{itemize}
    \item Test the Conductor mode in complete performances and tune the
    acceleration threshold, minimum speed, refractory period, and melodic
    movement distance.
    \item Measure the delay between the physical downstroke and the audible
    onset to identify whether the camera, MediaPipe, MIDI output, or synthesizer
    contributes most to the response time.
    \item Improve visual feedback by showing when the right-hand stroke is
    accepted and which chord degree will be applied at the next beat.
    \item Evaluate the reliability of all seven right-hand finger patterns,
    especially the thumb gestures, and adjust the gesture filtering if needed.
    \item Record short user trials to determine whether the separation between
    right-hand conducting and left-hand melody feels natural and musically
    satisfying.
\end{itemize}

\section{Resources}

\begin{itemize}
    \item Project repository:
    \url{https://github.com/YorgoAzzi/MarkovModel}
    \item MediaPipe Hand Landmarker documentation:
    \url{https://ai.google.dev/edge/mediapipe/solutions/vision/hand_landmarker}
    \item OpenCV documentation:
    \url{https://docs.opencv.org/}
\end{itemize}

\end{document}
